\begin{abstract}
LLM inference increasingly dominates the operational energy of deployed AI
services, yet energy per token changes substantially with workload shape,
serving configuration, hardware, and scale.  Exhaustively profiling every
candidate on its target GPU cluster is impractical when accelerator
allocations are scarce.  We present \system{}, a bounded hybrid agent that
treats tuning as an evidence-allocation problem: which configuration should be
examined next, at what fidelity, and when is real-GPU verification worth its
cost?  A constrained LLM planner maps natural-language intent to typed
subgoals, while deterministic guards select candidate--fidelity actions,
enforce GPU-hour and SLO constraints, maintain the Pareto state, and permit
recommendations only after target-level verification.  Its evidence substrate,
\sandbox{}, combines a CPU-resident projector, sparse measured correction,
restricted GPU execution, immutable provenance, and explicit abstention.  On
one H100 80GB serving Qwen2.5-7B with vLLM, we calibrate workload transfer under
one fixed configuration and evaluate the frozen model on two preregistered
holdouts totaling \BlindMeasurements{} real runs.  Energy MAPE is
\CalibratedEnergyMAPE{} on an eight-workload holdout and
\ConfirmationEnergyMAPE{} on a disjoint nine-workload confirmation.  A
preregistered gate supports TTFT at concurrency four
(\KFourTTFTMAPE{} MAPE) and abstains below four (\SparseTTFTMAPE{}), showing
that cheap evidence can guide decisions without being mislabeled as hardware
truth.  The resulting system provides a reproducible foundation for
budget-aware agentic configuration search on HPC clusters.
\end{abstract}

\begin{IEEEkeywords}
agentic AI, LLM inference, energy efficiency, multi-fidelity optimization,
Pareto search, HPC clusters
\end{IEEEkeywords}
